% !TEX TS-program = pdflatex
% !TEX encoding = UTF-8 Unicode

% This is a simple template for a LaTeX document using the "article" class.
% See "book", "report", "letter" for other types of document.

\documentclass[11pt]{article} % use larger type; default would be 10pt

\usepackage[utf8]{inputenc} % set input encoding (not needed with XeLaTeX)

%%% Examples of Article customizations
% These packages are optional, depending whether you want the features they provide.
% See the LaTeX Companion or other references for full information.

%%% PAGE DIMENSIONS
\usepackage{geometry} % to change the page dimensions
\geometry{a4paper} % or letterpaper (US) or a5paper or....
% \geometry{margin=2in} % for example, change the margins to 2 inches all round
% \geometry{landscape} % set up the page for landscape
%   read geometry.pdf for detailed page layout information

\usepackage{graphicx} % support the \includegraphics command and options

% \usepackage[parfill]{parskip} % Activate to begin paragraphs with an empty line rather than an indent

%%% PACKAGES
\usepackage{booktabs} % for much better looking tables
\usepackage{array} % for better arrays (eg matrices) in maths
\usepackage{paralist} % very flexible & customisable lists (eg. enumerate/itemize, etc.)
\usepackage{verbatim} % adds environment for commenting out blocks of text & for better verbatim
\usepackage{subfig} % make it possible to include more than one captioned figure/table in a single float
% These packages are all incorporated in the memoir class to one degree or another...

%%% HEADERS & FOOTERS
\usepackage{fancyhdr} % This should be set AFTER setting up the page geometry
\pagestyle{fancy} % options: empty , plain , fancy
\renewcommand{\headrulewidth}{0pt} % customise the layout...
\lhead{}\chead{}\rhead{}
\lfoot{}\cfoot{\thepage}\rfoot{}

%%% SECTION TITLE APPEARANCE
\usepackage{sectsty}
\allsectionsfont{\sffamily\mdseries\upshape} % (See the fntguide.pdf for font help)
% (This matches ConTeXt defaults)

%%% ToC (table of contents) APPEARANCE
\usepackage[nottoc,notlof,notlot]{tocbibind} % Put the bibliography in the ToC
\usepackage[titles,subfigure]{tocloft} % Alter the style of the Table of Contents
\renewcommand{\cftsecfont}{\rmfamily\mdseries\upshape}
\renewcommand{\cftsecpagefont}{\rmfamily\mdseries\upshape} % No bold!

\usepackage{amsmath}


%%% END Article customizations

%%% The "real" document content comes below...

\title{Examen 3}
\author{Profe Hugo y Gemini assistant}
%\date{} % Activate to display a given date or no date (if empty),
         % otherwise the current date is printed 

\begin{document}
\maketitle

 Guía de solución paso a paso (para el docente) - este tipo de tasas de cambio son la base para calcular latencias dinámicas en redes móviles (MANETs) o algoritmos de enrutamiento dinámico.

Parte a) Rapidez de expansión del radio de señal
Identificar variables y tasas conocidas:


Área del círculo: $A = \pi r^2$


Tasa de cambio del área: $\frac{dA}{dt} = 120\text{ m}^2/\text{s}$


Instante de evaluación: $A = 400\pi\text{ m}^2$


Incógnita: $\frac{dr}{dt}$


Encontrar el radio en el instante dado:


$$400\pi = \pi r^2 \implies r^2 = 400 \implies r = 20\text{ m}$$
Derivar implícitamente respecto al tiempo $t$:


$$\frac{dA}{dt} = 2\pi r \frac{dr}{dt}$$
Sustituir y despejar $\frac{dr}{dt}$:


$$120 = 2\pi (20) \frac{dr}{dt}$$
$$120 = 40\pi \frac{dr}{dt}$$
$$\frac{dr}{dt} = \frac{120}{40\pi} = \frac{3}{\pi} \approx 0.955\text{ m/s}$$
Parte b) Razón de cambio de la distancia entre los drones
Plantear el modelo geométrico: Las trayectorias hacia el este ($x$) y hacia el norte ($y$) forman un triángulo rectángulo con el servidor en el vértice del ángulo recto. Por el teorema de Pitágoras:


$$D^2 = x^2 + y^2$$
Identificar tasas y valores en $t = 10\text{ s}$:


$\frac{dx}{dt} = 4\text{ m/s} \implies \text{En } t = 10\text{ s}, \quad x = 4 \cdot 10 = 40\text{ m}$


$\frac{dy}{dt} = 3\text{ m/s} \implies \text{En } t = 10\text{ s}, \quad y = 3 \cdot 10 = 30\text{ m}$


Distancia $D$ en $t = 10\text{ s}$:


$$D = \sqrt{40^2 + 30^2} = \sqrt{1600 + 900} = \sqrt{2500} = 50\text{ m}$$
Derivar implícitamente respecto al tiempo $t$:


$$2D \frac{dD}{dt} = 2x \frac{dx}{dt} + 2y \frac{dy}{dt}$$
Dividiendo toda la ecuación entre 2:


$$D \frac{dD}{dt} = x \frac{dx}{dt} + y \frac{dy}{dt}$$
Sustituir y despejar $\frac{dD}{dt}$:


$$50 \frac{dD}{dt} = (40)(4) + (30)(3)$$
$$50 \frac{dD}{dt} = 160 + 90$$
$$50 \frac{dD}{dt} = 250$$
$$\frac{dD}{dt} = \frac{250}{50} = 5\text{ m/s}$$
¿Por qué funciona bien este ejercicio para primer año?
Enfoque gradual: La parte a) calienta motores con una sola variable dependiente ($r$), ideal para que pierdan el miedo a la regla de la cadena aplicada al tiempo. La parte b) introduce el clásico triángulo rectángulo (con una terna pitagórica $3-4-5$ que evita raíces cuadradas complicadas y permite que se concentren en el cálculo y no en la aritmética).


Conexión con la carrera: Puedes mencionar en clase que 

¿Te gustaría que agregue una tercera parte más desafiante (por ejemplo, involucrando el ángulo de visión entre los drones) o prefieres otro enfoque temático, como procesamiento de datos en servidores?



\end{document}
